\documentclass[11pt]{article}
% ===== Shared preamble for Calendar notes =====
\usepackage[margin=1in]{geometry}
\usepackage{amsmath,amssymb}
\usepackage{enumitem}
\usepackage{xcolor}
\usepackage{tcolorbox}
\tcbuselibrary{skins,breakable}
\usepackage{titlesec}
\usepackage{fancyhdr}
\usepackage{hyperref}
\usepackage{array}
\usepackage{booktabs}
\usepackage{longtable}
\usepackage{pifont} % for \ding

% ---- Colors ----
\definecolor{primary}{HTML}{1B3A6B}
\definecolor{accent}{HTML}{2E7D32}
\definecolor{warnclr}{HTML}{B45309}
\definecolor{tipclr}{HTML}{0F5D8C}
\definecolor{lightbg}{HTML}{F4F6F8}

% ---- Section styling ----
\titleformat{\section}
  {\normalfont\Large\bfseries\color{primary}}
  {\thesection}{0.5em}{}[{\color{primary}\titlerule[1pt]}]
\titleformat{\subsection}
  {\normalfont\large\bfseries\color{primary}}
  {\thesubsection}{0.5em}{}

% ---- Header/Footer ----
\pagestyle{fancy}
\fancyhf{}
\renewcommand{\headrulewidth}{0.4pt}
\fancyfoot[C]{\thepage}

% ---- Hyperref ----
\hypersetup{
  colorlinks=true,
  linkcolor=primary,
  urlcolor=tipclr,
  citecolor=accent
}

% ---- Custom boxes ----
% Formula / Key concept box
\newtcolorbox{formulabox}[1][]{
  enhanced, breakable,
  colback=lightbg, colframe=primary,
  boxrule=0.8pt, arc=2pt,
  left=8pt, right=8pt, top=6pt, bottom=6pt,
  title={\textbf{#1}},
  fonttitle=\bfseries\color{white},
  colbacktitle=primary,
  attach boxed title to top left={xshift=8pt, yshift=-2mm},
  boxed title style={boxrule=0pt, arc=2pt}
}

% Tip / trick box
\newtcolorbox{tipbox}[1][Tip]{
  enhanced, breakable,
  colback=blue!4, colframe=tipclr,
  boxrule=0.8pt, arc=2pt,
  left=8pt, right=8pt, top=6pt, bottom=6pt,
  title={\textbf{\textcolor{white}{\ding{72}\ #1}}},
  fonttitle=\bfseries, colbacktitle=tipclr,
}

% Warning / common mistake box
\newtcolorbox{warnbox}[1][Watch Out]{
  enhanced, breakable,
  colback=orange!6, colframe=warnclr,
  boxrule=0.8pt, arc=2pt,
  left=8pt, right=8pt, top=6pt, bottom=6pt,
  title={\textbf{\textcolor{white}{#1}}},
  fonttitle=\bfseries, colbacktitle=warnclr,
}

% Example / worked question box
\newtcolorbox{examplebox}[1][Example]{
  enhanced, breakable,
  colback=green!4, colframe=accent,
  boxrule=0.8pt, arc=2pt,
  left=8pt, right=8pt, top=6pt, bottom=6pt,
  title={\textbf{\textcolor{white}{#1}}},
  fonttitle=\bfseries, colbacktitle=accent,
}


\title{\color{primary}\textbf{Calendar --- Question Pattern Catalog}\\[4pt]\large Every Way Placement Exams Phrase a Calendar Question, Grouped by the Trick That Solves It}
\author{Aptitude Grind}
\date{}

\lhead{\small Calendar --- Question Pattern Catalog}
\rhead{\small Aptitude Grind}

\begin{document}
\maketitle
\thispagestyle{fancy}

Calendar questions look varied, but almost every one of them reduces to a handful of tricks from
\texttt{02-fast-tricks-and-shortcuts.tex}. This file catalogues the actual \textbf{phrasings} you'll
meet in TCS NQT, Infosys, Wipro, Capgemini, and CAT-style tests, grouped by category, each with the
fastest method and a verified mini-example. Skim this before a test to pattern-match a question to
its solution method in seconds.

\section{Category A --- Direct Day Lookup}

\begin{examplebox}[A1. Find the Day of a Given Date]
\emph{``What day of the week was 15 August 1947?''} \\
The most common calendar question by far --- direct application of the odd-days / anchor method.
\textbf{Method:} \S1 of Notes 02, or the mod-400 method in Notes 01. \\
\textbf{Answer:} 15 August 1947 was a \textbf{Friday}.
\end{examplebox}

\section{Category B --- Relative Day Arithmetic (No Calendar Needed)}

\begin{examplebox}[B1. Day After N Days]
\emph{``Today is Tuesday. What day will it be after 100 days?''} \\
\textbf{Method:} $N \bmod 7$, step forward. $100 \bmod 7 = 2 \to$ Tuesday $+2 = $ \textbf{Thursday}.
\end{examplebox}

\begin{examplebox}[B2. Day Before N Days]
\emph{``100 days before Friday is which day?''} \\
\textbf{Method:} $N \bmod 7$, step backward. $100 \bmod 7 = 2 \to$ Friday $-2 = $ \textbf{Wednesday}.
\end{examplebox}

\begin{examplebox}[B3. Yesterday / Tomorrow Questions]
\emph{``Tomorrow is Monday. What was yesterday?''} \\
\textbf{Method:} Anchor on the one day you're told, count off by $\pm 1$. Tomorrow = Monday
$\Rightarrow$ Today = Sunday $\Rightarrow$ Yesterday = \textbf{Saturday}.
\end{examplebox}

\begin{examplebox}[B4. Chained Relative-Day Questions]
\emph{``Three days after tomorrow is Friday. What day is today?''} \\
\textbf{Method:} Work backward. ``Three days after tomorrow'' = today $+4$ = Friday
$\Rightarrow$ today = Friday $-4$ = \textbf{Monday}.
\end{examplebox}

\begin{examplebox}[B5. Mirror-Day Questions]
\emph{``If yesterday was Thursday, what day will it be 100 days after tomorrow?''} \\
\textbf{Method:} Pin down today first (yesterday = Thursday $\Rightarrow$ today = Friday), then
tomorrow = Saturday, then apply B1: $100 \bmod 7 = 2 \to$ Saturday $+2 = $ \textbf{Monday}.
\end{examplebox}

\begin{examplebox}[B6. Calendar + Clock Combined]
\emph{``It is Monday, 11 PM. What day and time will it be after 30 hours?''} \\
\textbf{Method:} Split into full days + leftover hours: $30 = 24+6$. One full day moves Monday
$\to$ Tuesday; the remaining 6 hours push 11 PM past midnight, landing on
\textbf{Wednesday, 5 AM}. Always check whether leftover hours cross midnight.
\end{examplebox}

\section{Category C --- Nth / Last Weekday of a Month}

\begin{examplebox}[C1. Nth Weekday of a Month]
\emph{``The third Monday of July 2026 falls on which date?''} \\
\textbf{Method:} Find the weekday of the 1st, add 7 per subsequent occurrence.
1 July 2026 = Wednesday $\to$ first Monday = 6th $\to$ third Monday $= 6+14 = $ \textbf{20th}.
\end{examplebox}

\begin{examplebox}[C2. Last Weekday of a Month]
\emph{``Find the last Friday of February 2026 (28-day, non-leap).''} \\
\textbf{Method:} Find the weekday of the last day, count back in steps of 7 to the target weekday.
\textbf{Answer: 27 February 2026} (verified directly against the month).
\end{examplebox}

\begin{examplebox}[C3. Which Weekdays Occur 5 Times]
\emph{``In a 31-day month starting on a Tuesday, which weekdays occur 5 times?''} \\
\textbf{Method:} \S5 rule of Notes 02 --- $L = \text{days}-28$. For 31 days, $L=3$: the starting
weekday and the next two. Starting Tuesday $\Rightarrow$ \textbf{Tuesday, Wednesday, Thursday}.
\end{examplebox}

\begin{examplebox}[C4. February-Specific Frequency]
\emph{``In a leap year, which weekday occurs 5 times in February?''} \\
\textbf{Method:} 29 days $\Rightarrow L=1 \Rightarrow$ only the weekday 1 February falls on. In a
non-leap February (28 days, $L=0$), \textbf{no} weekday occurs 5 times --- all seven occur exactly
4 times.
\end{examplebox}

\section{Category D --- Same Calendar / Repeating Year}

\begin{examplebox}[D1. Same Calendar as a Given Year]
\emph{``The calendar for 2026 will be the same as the calendar for which future year?''} \\
\textbf{Method:} Running-odd-day-total from Notes 02 \S7 --- don't assume a fixed gap.
\textbf{Answer for 2026 specifically: 2037} (verified; 2026 is ordinary, gap = 11 years).
\end{examplebox}

\begin{examplebox}[D2. Odd-One-Out Calendar]
\emph{``Which of these years has a calendar different from the others: 1997, 2003, 2014, 2025?''} \\
\textbf{Method:} Two years share a calendar only if both are the same type (ordinary/leap)
\textbf{and} 1 January falls on the same weekday. Check leap/non-leap status first to eliminate
options fast, then compare 1 January weekdays only among years of matching type.
\end{examplebox}

\begin{examplebox}[D3. Calendar Repeats After N Years (Cycle Length)]
\emph{``After how many years does a leap year's calendar typically repeat?''} \\
\textbf{Method:} Table in Notes 02 \S7 --- typically \textbf{28 years} for a leap year (rarely 40,
only near a century that breaks the 4-year leap rule). For an ordinary year, typically
\textbf{6 or 11 years} (rarely 12). Never assume 17 --- that number doesn't actually occur.
\end{examplebox}

\section{Category E --- Leap Year Reasoning}

\begin{examplebox}[E1. Is This Year a Leap Year?]
\emph{``Is 2100 a leap year?''} \\
\textbf{Method:} Notes 02 \S4 quick check. 2100 is divisible by 4 and by 100, but \textbf{not} by
400 $\Rightarrow$ \textbf{not a leap year}. Classic trap --- most people stop at ``divisible by 4.''
\end{examplebox}

\begin{examplebox}[E2. Counting Leap Years in a Range]
\emph{``How many leap years are there from 1601 to 2000?''} \\
\textbf{Method:} Notes 02 \S4 formula, $\text{count}(2000)-\text{count}(1600) = 485-388 = $
\textbf{97}.
\end{examplebox}

\begin{examplebox}[E3. Century / Leap-Year Reasoning (CAT-style)]
\emph{``Can the last day of a century (a year ending in 00) ever fall on a Tuesday?''} \\
\textbf{Method:} Centuries only ever contribute 5, 3, 1, or 0 odd days (for 100, 200, 300, 400-year
blocks) --- never 2, 4, or 6. So a century's last day can only be \textbf{Sunday, Monday, Wednesday,
or Friday}, never Tuesday, Thursday, or Saturday.
\end{examplebox}

\section{Category F --- Date-to-Date and Span Problems}

\begin{examplebox}[F1. Days Between Two Dates]
\emph{``How many days are there between 15 August 2010 and 25 December 2010?''} \\
\textbf{Method:} Convert both to a running day-of-year count (or subtract directly), being careful
whether inclusive or exclusive is wanted. \textbf{Answer: 132 days} (exclusive of the start date).
\end{examplebox}

\begin{examplebox}[F2. Consecutive Dates Within a Month]
\emph{``If 18 January is a Tuesday, what day is 27 January?''} \\
\textbf{Method:} Notes 02 \S6 rule --- difference in days mod 7. $27-18=9$, $9 \bmod 7 = 2 \to$
Tuesday $+2 = $ \textbf{Thursday}.
\end{examplebox}

\begin{examplebox}[F3. Same Date, Different Year (Leap-Day Trap)]
\emph{``If 15 March 2015 was a Sunday, what day was 15 March 2016?''} \\
\textbf{Method:} Don't just add 1 odd day for ``one year passing'' --- check whether 29 February
falls \textbf{inside the specific span} being bridged. Here 29 Feb 2016 falls between the two dates,
so the gap is \textbf{2} odd days, not 1. Sunday $+2 \to$ \textbf{Tuesday}.
\end{examplebox}

\section{Category G --- Reverse Problems (Given a Clue, Find the Date/Day)}

\begin{examplebox}[G1. Given the Nth Occurrence, Find the 1st Occurrence]
\emph{``The third Saturday of a month falls on the 16th. Find the date of the first Saturday.''} \\
\textbf{Method:} Subtract $7\times(N-1)$ days. Third minus 2 weeks: $16-14=2$. \textbf{Answer: 2nd.}
\end{examplebox}

\begin{examplebox}[G2. Given the Month's Starting Day, Find a Later Date's Day]
\emph{``A month starts on Wednesday and has 30 days. What day is the 30th?''} \\
\textbf{Method:} Notes 02 \S6 rule --- $(30-1) \bmod 7 = 29 \bmod 7 = 1 \to$ Wednesday $+1 = $
\textbf{Thursday}.
\end{examplebox}

\begin{examplebox}[G3. Calendar Matching --- Find All Occurrences of a Weekday]
\emph{``If the 15th of a month is a Friday, what other dates in that month are Fridays?''} \\
\textbf{Method:} Every matching weekday is exactly 7 days apart. From the 15th: 1, 8,
\textbf{15, 22, 29} (list depends on month length --- always check both directions from the known
date, not just forward).
\end{examplebox}

\section{Category H --- Logic Puzzles and Word Problems}

\begin{examplebox}[H1. Birthday / Recurrence Puzzles]
\emph{``A person was born on a Monday. Is it possible for their birthday to fall on a Monday again
every single year?''} \\
\textbf{Method:} A birthday repeats on the same weekday only if the gap between birthdays
contributes 0 odd days --- impossible every year, since every ordinary year contributes 1 odd day
and every leap year contributes 2. It can only repeat on the same weekday in the rare years where
the running odd-day total (Notes 02 \S7) resets to a multiple of 7, not annually.
\end{examplebox}

\begin{examplebox}[H2. Working-Day Counting]
\emph{``An office is closed every Sunday. If a 31-day month starts on a Saturday, how many working
days does it have?''} \\
\textbf{Method:} Count Sundays using the weekday-frequency rule (\S C3/C4), subtract from total
days. A 31-day month starting on Saturday has $L=3$ extra weekdays (Sat, Sun, Mon occur 5 times) ---
so Sunday occurs \textbf{5} times $\Rightarrow$ working days $= 31-5 = $ \textbf{26}.
\end{examplebox}

\begin{examplebox}[H3. Multi-Clue Logic Puzzles]
\emph{``1 January of a certain year is a Friday, and 1 January of the following year is a Sunday.
What can you conclude about the certain year?''} \\
\textbf{Method:} The shift from one 1 January to the next equals that year's odd-day contribution.
Friday $\to$ Sunday is a shift of \textbf{2} odd days, which only a \textbf{leap year} contributes
$\Rightarrow$ the certain year must be a leap year.
\end{examplebox}

\section{Priority Order for Limited Prep Time}

If you're short on time, drill these first --- they cover the large majority of what actually
appears in placement tests, roughly in order of frequency:

\begin{enumerate}[leftmargin=1.6em]
  \item \textbf{A1} --- Find the day for a given date (mod-400 or anchor method)
  \item \textbf{B1/B2} --- Day after/before N days (pure $\bmod 7$)
  \item \textbf{E1} --- Leap year identification (the century exception is the \#1 trap)
  \item \textbf{C1} --- Nth weekday of a month
  \item \textbf{D1} --- Same/repeating calendar year
  \item \textbf{C3/C4} --- Weekday frequency in 28/29/30/31-day months
  \item \textbf{F1} --- Days between two dates
  \item \textbf{F3} --- Same date, different year (leap-day-crossing trap)
  \item \textbf{E3} --- Century reasoning (CAT-style, less common but high-value when it appears)
\end{enumerate}

\end{document}
